%% kapa: exam=1 round=34 year=2023 period=1 form=A subject=civil_law subject_ko=민법 seq=27 number=27
%% source: Q-Net 감정평가사 기출문제 (cst00309) artlSeq=5211789 file=raw/1008/5211789/2023년도 제34회 감정평가사 1차시험 1교시.pdf
%% conversion: kordoc 4.24.51 -> tools/gen1.py (auto); review=proofread-sonnet; flags=none
\begin{question}{27}
점유자와 회복자의 관계에 관한 설명으로 옳은 것은? (다툼이 있으면 판례에 따름)
\begin{choices}
  \item 선의의 점유자라도 점유물의 사용으로 인한 이익은 회복자에게 반환하여야 한다.
  \item 임차인이 지출한 유익비는 임대인이 아닌 점유회복자에 대해서도 민법 제203조 제2항에 근거하여 상환을 청구할 수 있다.
  \item 과실수취권 있는 선의의 점유자란 과실수취권을 포함하는 본권을 가진다고 오신할 만한 정당한 근거가 있는 점유자를 가리킨다.
  \item 선의점유자에 대해서는 점유에 있어서의 과실(過失) 유무를 불문하고 불법행위를 이유로 한 손해배상책임이 배제된다.
  \item 점유물이 타주점유자의 책임 있는 사유로 멸실된 경우, 그가 선의의 점유자라면 현존이익의 범위에서 손해배상책임을 진다.
\end{choices}
\end{question}
\answer{3}
